\section{未来风险：能力增长之外的代价}

\subsection{从生产力到社会成本}

模型越来越强，并不意味着所有结果都朝正方向走。课程最后把视角推向了更现实的层面：AI 系统的收益和成本是同时扩张的。

\begin{figure}[H]
\centering
\includegraphics[width=0.95\textwidth]{slides-images/slide-055.jpg}
\caption{AI 发展总伴随着隐藏成本：错误信息、偏见、能源消耗、数据掠夺和劳动问题都是真实代价\protect\footnotemark}
\end{figure}
\footnotetext{Slides 直接列出了多种 human and environmental costs。}

这些成本并不是“未来可能出现”的抽象概念，而是已经存在的现实问题：
\begin{itemize}
    \item 虚假信息和误导内容；
    \item 偏见、幻觉和歧视性输出；
    \item 隐私泄露和版权争议；
    \item 数据采集、标注和算力背后的劳动与环境成本。
\end{itemize}

\begin{warningbox}{“更强”不等于“更安全”}
能力上升会同时放大正效应和负效应。系统越强，越需要在评测、部署、治理和责任边界上保持谨慎。
\end{warningbox}

\subsection{关于存在风险的讨论}

课尾还触及了 existential risk / extinctions risk 这类更激烈的话题。Manning 的处理方式并不是耸动，而是承认它是一个需要讨论的真实议题，同时也指出这不应遮蔽当下更具体、更可验证的风险。

\begin{itemize}
    \item 有些讨论强调极端远期风险；
    \item 另一些讨论则提醒我们不要忽视当前已经发生的系统性伤害；
    \item 对研究者来说，最重要的是把抽象担忧落到可观测的机制和政策上。
\end{itemize}

\subsection{本章小结}

AI 风险不是附录，而是模型能力增长的一部分。最后一讲把这件事说得很明确：如果只看 benchmark 上升，就会错过技术扩张带来的社会后果。

%% ========== 研究议程 ========== 

